This proof was given to me (Hugo Moeneclaey) by Thierry Coquand.

\subsection{Standard \'etale scheme}

\begin{definition}
Given $P,Q:R[X]$ with $P$ monic non-constant we define $\Std(P,Q) = \Spec((R[X]/P)_Q)$
\end{definition}

\begin{definition}
Assume given $P,Q:R[X]$ with $P$ monic non-constant.

We write $\psi(P,Q)$ for the proposition ``$\Std(P,Q)$ is formally \'etale''.

We write $\delta(P,Q)$ for the proposition $\neg\neg\Std(P,Q)$.
\end{definition}

\begin{definition}
Assume given $P:R[X]$ monic non-constant, we say that $P$ is unramifiable if $\delta(P,P')$.
\end{definition}

\begin{definition}
A standard \'etale scheme is a scheme merely of the form $\Std(P,Q)$ for $P$ monic non-constant such that $\psi(P,Q)$ holds.
\end{definition}

\begin{remark}\label{remark-unramifiable-polynomial}
Assume given $P:R[X]$ monic non-constant, then we have $\psi(P,P')$. So if $P$ is unramifiable, then $\Std(P,P')$ is in the \'etale topology.
\end{remark}

\begin{lemma}\label{standard-etale-characterisation}
Assume given $P,Q:R[X]$ with $P$ monic non-constant. Then $\psi(P,Q)$ if and only if $Q\in\sqrt((P,P'))$
\end{lemma}

\begin{proof}
TODO
\end{proof}

\begin{lemma}\label{delta-is-open}
Assume given $P,Q:R[X]$ with $P$ monic non-constant. Then $\delta(P,Q)$ is open.
\end{lemma}

\begin{proof}
  Following \cite{lombardi-quitte}[III.4], let \(\Adu_{R,P}\) be the universal splitting algebra and let \(x_1,\dots,x_n:\Adu_{R,P}\) be the roots of \(P\).
  Then \(D_{P,Q}\colonequiv Q(x_1)\neq 0 \vee\dots\vee Q(x_n)\neq 0\) is open since \(\Adu_{R,P}\) is finite free as an \(R\)-module (\cite{lombardi-quitte}[III.4.4]).
  
  To show that \(\delta(R,Q)\) is open, we will show that it is equivalent to \(D_{P,Q}\).
  We first show \(\neg\neg\Std(P,Q)\to D_{P,Q}\).
  By being open, \(D_{P,Q}\) is double negation stable, so we can assume \(x:\Std(P,Q)\) and show \(D_{P,Q}\). We get a homomorphism \(\Adu_{P,Q}\to (R[X]/P)_Q\) by sending all the \(x_i\) to \(x\).
  Assume \(Q(x_1)=0\wedge\dots\wedge Q(x_n)=0\), then we have \(Q(x)=0\), which is a contradiction.

  We conclude by showing \(\neg\Std(P,Q) \to \neg D_{P,Q}\).
  So assume \(\Std(P,Q)=\emptyset\) and therefore that for \(x:R[X]/P\), \(Q(x)\) is nilpotent.
  For any \(x_i:\Adu_{P,Q}\) we can look at the map \(R[X]/P\to \Adu_{P,Q}\) given by mapping \(x\) to \(x_i\). So any \(Q(x_i)\) is nilpotent and we get \(\neg D_{P,Q}\).
\end{proof}




\subsection{The local form for \'etale scheme}

\begin{theorem}
Assume given an \'etale scheme $X$. Then there exists a finite Zariski cover of $X$ by standard \'etale schemes. 
\end{theorem}

We give no proof, see Thierry's note at \url{https://arxiv.org/abs/2601.04143} for a constructive proof, although not in SAG.


\subsection{Characterisation for the \'etale sheaves using monic polynomials}

\begin{definition}
The \'etale monic sheafification is the lex modality localising $\Spec(R[X]/P)$ for all unramifiable monic $P$.
\end{definition}

So saying something holds under the \'etale monic sheafifcation says it holds assuming all monic unramifiable polynomials have roots. In this section we will prove that \'etale sheafification and \'etale monic sheafification coincide. We first give two auxiliary lemmas.

\begin{lemma}
Assume given $P,Q:R[X]$ with $P$ monic non-constant such that $\psi(P,Q)$ holds. Given $x:R$ such that $P(x)=0$, then:
\[\delta(P,Q) \Rightarrow Q(x)\not=0 \lor \delta(P/(X-x),Q)\]
\end{lemma}

\begin{proof}
Assume $\delta(P,Q)$. By \Cref{delta-is-open} we know that the goal is $\neg\neg$-stable, so we can assume $x:\Std(P,Q)$, i.e.\ $y:R$ such that $P(y)=0$ and $Q(y)\not=0$. We can also use $Q(x)=0$ or $Q(x)\not=0$. If $Q(x)\not=0$ we are done, assume $Q(x)=0$. Then we want to prove that $P/(X-x)(y) = 0$, but since $P(y)=0$ and $x\not=y$ this is clear.
\end{proof}

\begin{lemma}
Assume given $P,Q:R[X]$ with $P$ monic non-constant, if $\psi(P,Q)$ and $\delta(P,Q)$ then $\delta(P,P')$
\end{lemma}

\begin{proof}
Assume $\psi(P,Q)$ and $\delta(P,Q)$, we have to prove $\delta(P,P')$. By \Cref{delta-is-open} \rednote{(and possibly \Cref{remark-unramifiable-polynomial}?)} we have that our goal is $\neg\neg$-stable, so we can assume $\Std(P,Q)$, i.e. we can assume $x:R$ such that $P(x)=0$ and $Q(x)\not=0$. Then $P'(x)\not=0$ is enough to conclude. 

But if $P'(x)=0$ then let's prove that for all $\epsilon:R$ such that $\epsilon^2=0$, we have $\epsilon=0$. Indeed $P(x+\epsilon) = 0$ and since $Q(x)\not=0$ and $\neg\neg(\epsilon=0)$, we have that $Q(x+\epsilon)\not=0$, so that $x+\epsilon$ is in $\Std(P,Q)$. Then since $\psi(P,Q)$ we get that $x=x+\epsilon$, so that $\epsilon=0$.
\end{proof}

Now we are ready to prove that the standard \'etale scheme that are in the topology are inhabited under the \'etale monic sheafification, i.e.\ they are inhabited assuming unramifiable polynomials have roots.

\begin{lemma}
Assume given $P,Q:R[X]$ with $P$ monic non-constant, if $\psi(P,Q)$ and $\delta(P,Q)$ then $\propTrunc{\Std(P,Q)}_{\mathrm{EtMon}}$.
\end{lemma}

\begin{proof}
TODO
\end{proof}

We can extend to arbitrary elements in the topology.

\begin{lemma}
Assume given $S$ in the \'etale topology (i.e.\ $S$ affine, formally \'etale and non-empty), then $\propTrunc{S}_{\mathrm{EtMon}}$
\end{lemma}

\begin{proof}
TODO
\end{proof}

\begin{corollary}
A type is an \'etale sheaf if and only if it is an \'etale monic sheaf.
\end{corollary}



